% ==========================================
% PHAN A - STATIC ROUTING CHO KHACH HANG D1
% ==========================================

\section{Phần A -- Static routing cho khách hàng D1}
\label{sec:partA}

\subsection{Mục tiêu}

Phần này cấu hình định tuyến tĩnh cho D1,
là router khách hàng không chạy giao thức định tuyến động.
D1 sử dụng default route qua C1 trên L10 làm đường chính,
đồng thời có default route qua B2 trên L11 với AD 200
làm đường dự phòng.
C1 và B2 được cấu hình static route về LAN-D
để có đường chuyển tiếp gói tin tới mạng khách hàng.

Theo bảng quy hoạch của nhóm K = 7126,
LAN-D là mạng \texttt{10.171.126.192/26},
được mô phỏng bằng Loopback1 của D1
với địa chỉ \texttt{10.171.126.193}.
Hai liên kết tới D1 gồm:
\begin{itemize}
    \item L10: C1 (\texttt{10.71.26.37})
    nối tới D1 (\texttt{10.71.26.38}).
    \item L11: B2 (\texttt{10.71.26.41})
    nối tới D1 (\texttt{10.71.26.42}).
\end{itemize}

\subsection{Triển khai cấu hình}

\subsubsection{Ý nghĩa các lệnh chính}

Static route được khai báo bằng lệnh:

\begin{lstlisting}
ip route <destination> <subnet-mask> <next-hop> [AD]
\end{lstlisting}

Trong đó, destination và subnet mask xác định mạng đích;
next-hop là địa chỉ router kế tiếp.
AD (Administrative Distance) biểu thị mức ưu tiên
giữa các tuyến tới cùng một prefix: tuyến có AD thấp hơn
được ưu tiên khi đủ điều kiện cài đặt vào bảng định tuyến.
Static route không chỉ định AD sử dụng giá trị mặc định là 1.

Default route sử dụng địa chỉ mạng và subnet mask
đều bằng \texttt{0.0.0.0}, tương ứng prefix
\texttt{0.0.0.0/0}. Router dùng tuyến này khi
không có tuyến cụ thể hơn phù hợp với địa chỉ đích.
Trên D1, tuyến qua C1 có AD 1,
còn tuyến qua B2 có AD 200.
Tuyến qua B2 là floating static route:
chỉ được chọn khi tuyến chính không còn được cài đặt
và đường tới next-hop dự phòng vẫn hợp lệ.

\subsubsection{Cấu hình trên D1}

Nhóm cấu hình hai default route trên D1 như sau:

\begin{lstlisting}
conf t
ip route 0.0.0.0 0.0.0.0 10.71.26.37
ip route 0.0.0.0 0.0.0.0 10.71.26.41 200
end
write memory
\end{lstlisting}

\subsubsection{Cấu hình đường về LAN-D trên C1 và B2}

Để gói tin từ phía ISP có thể tới LAN-D,
C1 và B2 cần static route trỏ về D1.
Mạng đích là \texttt{10.171.126.192/26},
có subnet mask \texttt{255.255.255.192}.

\textbf{Trên C1:}

\begin{lstlisting}
conf t
ip route 10.171.126.192 255.255.255.192 10.71.26.38
end
write memory
\end{lstlisting}

\textbf{Trên B2:}

\begin{lstlisting}
conf t
ip route 10.171.126.192 255.255.255.192 10.71.26.42
end
write memory
\end{lstlisting}

Hai static route này cung cấp đường về mạng khách hàng
và làm cơ sở để quảng bá LAN-D bằng BGP ở Phần D.

\subsection{Kiểm chứng kết quả}

\subsubsection{Kiểm tra hai default route trong cấu hình D1}

Trên D1, nhóm kiểm tra các dòng cấu hình định tuyến tĩnh:

\begin{lstlisting}
show running-config | include ip route
\end{lstlisting}

\partzeroimage
{images/configuration/part_A/D1_static.png}
{Cấu hình default route chính và dự phòng trên D1-7126}
{fig:partA-d1-static}

Hình~\ref{fig:partA-d1-static} xác nhận D1 có đủ
hai default route: tuyến qua C1 dùng AD mặc định là 1,
còn tuyến qua B2 được chỉ định AD 200.

\subsubsection{Kiểm tra bảng định tuyến trên D1}

\begin{lstlisting}
show ip route
\end{lstlisting}

\partzeroimage
{images/configuration/part_A/D1_route.png}
{Bảng định tuyến trên D1-7126 khi đường chính qua C1 đang hoạt động}
{fig:partA-d1-route}

Hình~\ref{fig:partA-d1-route} ghi nhận
\texttt{Gateway of last resort} là \texttt{10.71.26.37}
và tuyến đang được sử dụng:

\begin{lstlisting}
S* 0.0.0.0/0 [1/0] via 10.71.26.37
\end{lstlisting}

Ký hiệu \texttt{S} cho biết đây là static route,
dấu \texttt{*} đánh dấu tuyến ứng viên làm default route.
Trong \texttt{[1/0]}, giá trị 1 là AD và 0 là metric.
Tuyến qua B2 có AD 200 vẫn nằm trong cấu hình,
nhưng chưa được chọn vì tuyến qua C1 có AD thấp hơn.

Phép thử chuyển sang default route dự phòng khi gây lỗi
thuộc nội dung F1.

\subsubsection{Kiểm tra static route về LAN-D trên C1 và B2}

Trên cả C1 và B2, nhóm sử dụng lệnh:

\begin{lstlisting}
show running-config | include ip route
\end{lstlisting}

\partzeroimage
{images/configuration/part_A/C1_static.png}
{Static route về LAN-D trên C1-7126 qua liên kết L10}
{fig:partA-c1-static}

Hình~\ref{fig:partA-c1-static} xác nhận C1 có tuyến tới
\texttt{10.171.126.192/26} qua next-hop
\texttt{10.71.26.38}, là địa chỉ của D1 trên L10.

\partzeroimage
{images/configuration/part_A/B2_static.png}
{Static route về LAN-D trên B2-7126 qua liên kết L11}
{fig:partA-b2-static}

Hình~\ref{fig:partA-b2-static} xác nhận B2 định tuyến
LAN-D /26 qua D1 \texttt{10.71.26.42} trên L11.

\newpage
